\subsubsection{Subproceso 3.B: Confirmación y Registro Transaccional Oficial}

\textbf{Inicio:} El subproceso se activa cuando el Subproceso 3.A concluye la verificación y control de reglas y emite la Pre-Ficha de Matrícula Aprobada o Subsanada, acompañada de los dictámenes de rectificación resueltos por la Dirección de Escuela. En ese momento la plataforma Intranet DIRCE cierra la fase de edición del estudiante y la carga académica queda lista para su asentamiento en el SIGA.

\textbf{Incluye:}
\begin{itemize}
    \item Asentamiento transaccional definitivo de las cargas lectivas aprobadas en la base de datos central del SIGA, con bloqueo de la fase de edición en la Intranet DIRCE.
    \item Generación del Reporte de Matrícula, documento informativo individual que acredita las asignaturas registradas para el ciclo en curso.
    \item Semana de depuración y Cierre de Matrícula por Facultad: consolidación del padrón depurado y remisión del archivo digital oficial de la Oficina de Estadística de la Facultad hacia la DIRCE (RR-0570-2022, Art. 78°).
    \item Emisión de la Boleta de Matrícula Oficial en el SIGA con visado institucional de la Oficina de Estadística.
    \item Generación y distribución de los padrones de clase a los docentes el primer día hábil de la segunda semana lectiva (RR-0570-2022, Art. 78°).
    \item Reporte institucional a SUNEDU y MINEDU dentro del plazo de veinte días (RR-0570-2022, Art. 79°) y habilitación del trámite de expedición del carné universitario.
\end{itemize}

\textbf{Fin:} El subproceso concluye cuando la DIRCE consolida el cierre de matrícula y, en un plazo no mayor a diez días hábiles, pone a disposición de las Facultades el registro de matrícula, las boletas y la plataforma del SIGA actualizada (RR-0570-2022, Art. 80°), con la remisión previa de los reportes oficiales a SUNEDU y MINEDU; con ello la matrícula del periodo queda oficializada.

\textbf{Fuera del alcance:}
\begin{itemize}
    \item La acreditación del estudiante, la verificación de compromisos financieros y la asignación del turno, que corresponden al Subproceso 1.
    \item La selección de asignaturas, secciones y horarios, y el registro de la Pre-Ficha, propios del Subproceso 2.
    \item La auditoría de prerrequisitos, la gestión del Grupo Cero y la resolución de solicitudes de rectificación, que se realizan en el Subproceso 3.A.
    \item Las modificaciones de matrícula posteriores al cierre, que solo proceden mediante resolución decanal o rectoral extemporánea y se tramitan fuera del flujo regular del subproceso.
\end{itemize}
